\chapter{Développement RAW et profil}
\section{Choisir le rendu de départ}
Le groupe \emph{Développement RAW / Profil} permet de choisir un profil DCP
compatible avec la photo. Les propositions dépendent du boîtier, des profils
disponibles et du build. Vérifier la description du profil effectivement utilisé,
pas seulement le nom demandé dans une liste.

Un profil appareil définit un rendu de départ : il peut changer les couleurs,
le contraste et la luminosité apparente avant les retouches. Deux logiciels
ne produisent pas nécessairement les mêmes pixels avec un intitulé similaire.
Ce groupe n’applique pas un DCP d’appareil à un JPEG comme s’il s’agissait d’un RAW.

\section{Changer de profil}
Choisir une entrée dans la liste, attendre le calcul puis comparer le rendu.
Le moteur peut éviter un redécodage complet lorsqu’un tampon compatible est
disponible ; l’interface indique cette capacité. Elle ne garantit pas que
tous les changements de développement sont instantanés.
Les commandes susceptibles de modifier le développement sont désactivées
pendant certains traitements.

\section{Balance automatique et calibration}
La case \emph{Balance des blancs automatique} se trouve ici, une seule fois.
Les réglages manuels de température et de nuance restent dans \emph{Couleur}.
Ne pas confondre balance automatique, température de prise de vue et profil DCP.

La calibration RAW est une option avancée des préférences, accessibles avec
\touche{Ctrl + virgule}, onglet \emph{Avancé · RAW}. Sur les boîtiers et modes
compatibles, le moteur dispose d’une calibration par défaut. Vérifier son état
appliqué dans la description : une case activée n’assure pas qu’un fichier
inconnu possède une caractérisation disponible.

\attention{Le choix avancé concerne l’image ouverte et son compagnon,
pas une nouvelle préférence globale pour toutes les prochaines photos.
Ne changer le manifeste de calibration que dans un protocole maîtrisé.}

\paragraph{Exemple.}
Pour comparer deux profils, garder les retouches identiques, désactiver les
corrections non nécessaires au test, changer uniquement le profil et examiner
les tons de peau, les couleurs vives et les hautes lumières. Noter la balance
des blancs choisie pour rendre la comparaison reproductible.
